Ao meu orientador, Prof. Dr. César Lincoln Cavalcante Mattos, pela orientação sempre atenta, pelo acompanhamento semanal durante toda a pesquisa e, em especial, pela paciência nas muitas vezes em que mudei de ideia ao longo do caminho.

Ao meu coorientador, Prof. Dr. João Paulo Pordeus Gomes, que foi meu primeiro contato na UFC e se mostrou dedicado desde o começo, a ponto de me explicar redes convolucionais em uma folha de papel logo na primeira vez que conversamos pessoalmente.

Aos professores que compuseram a banca examinadora, Prof. Dr. José Antônio Fernandes de Macêdo, de quem também tive a satisfação de ser aluno e que tornou a ciência de dados mais interessante com tantas aplicações reais, e Prof. Dr. Diego Parente Paiva Mesquita, pela disponibilidade e pelas contribuições que enriqueceram este trabalho.

Ao Prof. Dr. Ronaldo Menezes, professor visitante no programa, pelo entusiasmo com que aproximou a turma de um tema novo e pela perspectiva diferente que trouxe de sua experiência no exterior.

Ao amigo e colega Manolidis Efstratios Júnior, que entrou no mestrado junto comigo e com quem dividi quase todas as disciplinas e as reuniões de orientação. Fazer esse percurso lado a lado tornou tudo mais leve.

Ao MDCC (Mestrado e Doutorado em Ciência da Computação) da Universidade Federal do Ceará, pelo ensino público e de qualidade e pela estrutura que tornou este mestrado possível.

Ao Doutorando em Engenharia Elétrica, Ednardo Moreira Rodrigues, e seu assistente, Alan Batista de Oliveira, aluno de graduação em Engenharia Elétrica, pela adequação do \textit{template} utilizado neste trabalho para que o mesmo ficasse de acordo com as normas da biblioteca da Universidade Federal do Ceará (UFC). %AVISO: Você pode usar este template uma vez que der os devidos créditos. Portanto, mantenha este parágrafo de agradecimento.
